\chapter{1991年上海卷}
\section{填空题}

\begin{problem}
函数 \(y = 3^x\)（\(x \geqslant 0\)）的反函数是\fillinblank{}.

\end{problem}

\begin{answer}
\(y = \log_3 x\)（\(x \geqslant 1\)）
\end{answer}

\begin{problem}
已知 \(\pi < \theta < \frac{3}{2}\pi\)，\(\cos \theta = -\frac{4}{5}\)，则 \(\cos \frac{\theta}{2} =\)\fillinblank{}.

\end{problem}

\begin{answer}
\(-\frac{\sqrt{10}}{10}\)
\end{answer}

\begin{problem}
一个圆柱的底面直径和高都等于一个球的直径，则这个圆柱的体积与球体积的比值为\fillinblank{}.

\end{problem}

\begin{answer}
\(\frac{3}{2}\)
\end{answer}

\begin{problem}
函数 \(y = \sin(\pi x + 2)\) 的最小正周期是\fillinblank{}.

\end{problem}

\begin{answer}
\(2\)
\end{answer}

\begin{problem}
椭圆 \(x^2 + 4y^2 = 1\) 的离心率是\fillinblank{}.

\end{problem}

\begin{answer}
\(\frac{\sqrt{3}}{2}\)
\end{answer}

\begin{problem}
在无穷等比数列 \(\{a_n\}\) 中，\(a_1 = 3\sqrt{3}\)，\(a_8 = \sqrt{3}\). 则 \(\lim\limits_{n \to \infty}(a_1 + a_3 + a_5 + \cdots + a_{2n-1}) =\)\fillinblank{}.

\end{problem}

\begin{answer}
\(\frac{9}{2}\sqrt{3}\)
\end{answer}

\begin{problem}
若母线长为 \(4\) 的圆锥的轴截面的面积为 \(8\)，则圆锥的侧面积为\fillinblank{}（结果中保留 \(\pi\)）.

\end{problem}

\begin{answer}
\(8\sqrt{2}\pi\)
\end{answer}

\begin{problem}
\(\left(2x^2 - \frac{1}{x}\right)^8\) 的展开式中的常数项为\fillinblank{}（结果用数值表示）.

\end{problem}

\begin{answer}
\(60\)
\end{answer}

\begin{problem}
设函数 \(y = f(x)\) 的图象关于直线 \(x = 1\) 对称，若当 \(x \leqslant 1\) 时，\(y = x^2 + 1\)，则当 \(x > 1\) 时，\(y = \)\fillinblank{}.

\end{problem}

\begin{answer}
\(x^2 - 4x + 5\)
\end{answer}

\begin{problem}
设有半径为 \(4\) 的圆，它在极坐标系内的圆心坐标是 \((4,\pi)\)，则这圆的极坐标方程是\fillinblank{}.

\end{problem}

\begin{answer}
\(\rho = -8\cos \theta\)
\end{answer}

\section{单选题}

\begin{problem}
若全集 \(I = \R\)，\(A = \{x \mid \sqrt{x + 1} \leqslant 0\}\)，\(B = \{x \mid \lg(x^2 - 2) = \lg x\}\)，则 \(A \cap \overline{B}\) 是（\quad）

\choices
    {\(\{2\}\)}
    {\(\{-1\}\)}
    {\(\{x \mid x \leqslant -1\}\)}
    {\(\varnothing\)}

\end{problem}

\begin{answer}
B
\end{answer}

\begin{problem}
在等差数列 \(\{a_n\}\) 中，若 \(a_3 + a_4 + a_5 + a_6 + a_7 = 450\)，则 \(a_2 + a_8\) 的值等于（\quad）

\choices
    {\(45\)}
    {\(75\)}
    {\(180\)}
    {\(300\)}

\end{problem}

\begin{answer}
C
\end{answer}

\begin{problem}
设长方体的对角线的长度是 \(4\)，过每一顶点有两条棱与对角线的夹角都是 \(60^{\circ}\)，则此长方体的体积是（\quad）

\choices
    {\(\frac{32}{72}\sqrt{3}\)}
    {\(8\sqrt{2}\)}
    {\(8\sqrt{3}\)}
    {\(16\sqrt{3}\)}

\end{problem}

\begin{answer}
B
\end{answer}

\begin{problem}
下列各函数图象中，表示函数 \(y = x^{-\frac{1}{2}}\) 的是（\quad）

\choices
    {\choicebitmap[width=0.15\paperwidth]{img/1991/shanghai/q14_choiceA.png}}
    {\choicebitmap[width=0.15\paperwidth]{img/1991/shanghai/q14_choiceB.png}}
    {\choicebitmap[width=0.15\paperwidth]{img/1991/shanghai/q14_choiceC.png}}
    {\choicebitmap[width=0.15\paperwidth]{img/1991/shanghai/q14_choiceD.png}}

\end{problem}

\begin{answer}
C
\end{answer}

\begin{problem}
下列四个命题中的假命题是（\quad）

\choices
    {存在这样的 \(\alpha\) 和 \(\beta\) 的值，使得 \(\cos(\alpha + \beta) = \cos \alpha \cos \beta + \sin \alpha \sin \beta\)}
    {不存在无穷多个 \(\alpha\) 和 \(\beta\) 的值，使得 \(\cos(\alpha + \beta) = \cos \alpha \cos \beta + \sin \alpha \sin \beta\)}
    {对于任意的 \(\alpha\) 和 \(\beta\)，\(\cos(\alpha + \beta) = \cos \alpha \cos \beta - \sin \alpha \sin \beta\)}
    {不存在这样的 \(\alpha\) 和 \(\beta\) 的值，使得 \(\cos(\alpha + \beta) \neq \cos \alpha \cos \beta - \sin \alpha \sin \beta\)}

\end{problem}

\begin{answer}
B
\end{answer}

\begin{problem}
设有编号为 \(1,2,3,4,5\) 的五个球和编号为 \(1,2,3,4,5\) 的五个盒子，现将这五个球投放入这五个盒内，要求每个盒内投放一个球，并且恰好有两个球的编号与盒子的编号相同，则这样的投放方法的总数为（\quad）

\choices
    {\(20\)}
    {\(30\)}
    {\(60\)}
    {\(120\)}

\end{problem}

\begin{answer}
A
\end{answer}

\begin{problem}
下列四个式子中，正确的是（\quad）

\choices
    {\(\sin\left(\arccos \frac{2}{3}\right) > \sin\left(\arccos \frac{1}{3}\right)\)}
    {\(\tan\left(\arccos \frac{2}{3}\right) > \tan\left(\arccos \frac{1}{3}\right)\)}
    {\(\sin\left[\arccos\left(-\frac{2}{3}\right)\right] > \sin\left[\arccos\left(-\frac{1}{3}\right)\right]\)}
    {\(\tan\left[\arccos\left(-\frac{2}{3}\right)\right] > \tan\left[\arccos\left(-\frac{1}{3}\right)\right]\)}

\end{problem}

\begin{answer}
D
\end{answer}

\begin{problem}
设直线 \(a\) 在平面 \(M\) 内，则平面 \(M\) 平行于平面 \(N\) 是直线 \(a\) 平行于平面 \(N\) 的（\quad）

\choices
    {充分条件但非必要条件}
    {必要条件但非充分条件}
    {充分必要条件}
    {非充分条件也非必要条件}

\end{problem}

\begin{answer}
A
\end{answer}

\begin{problem}
下列参数方程（\(t\) 为参数）与普通方程 \(x^2 - y = 0\) 表示同一曲线的方程是（\quad）

\choices
    {\(\left\{\begin{array}{l} x = |t|, y = t^2, \end{array}\right.\)}
    {\(\left\{\begin{array}{l} x = \cos t, y = \cos^2 t, \end{array}\right.\)}
    {\(\left\{\begin{array}{l} x = \tan t, y = \frac{1 + \cos 2t}{1 - \cos 2t}, \end{array}\right.\)}
    {\(\left\{\begin{array}{l} x = \tan t, y = \frac{1 - \cos 2t}{1 + \cos 2t}, \end{array}\right.\)}

\end{problem}

\begin{answer}
D
\end{answer}

\begin{problem}
在平面直角坐标系 \(xOy\) 中，设曲线 \(4x^3 - y^3 = 4\) 的一条倾角为锐角的渐近线为 \(L\)，现在平移坐标轴，把原点 \(O\) 移到 \(O'(3,-4)\)，则渐近线 \(L\) 在新坐标系 \(x'O'y'\) 中的方程是（\quad）

\choices
    {\(2x' - y' + 10 = 0\)}
    {\(2x' - y' - 10 = 0\)}
    {\(2x' + y' + 2 = 0\)}
    {\(2x' + y' - 2 = 0\)}

\end{problem}

\begin{answer}
A
\end{answer}

\section{解答题}

\begin{problem}
已知 \(\sin\left(x - \frac{3}{4}\pi\right)\cos\left(x - \frac{\pi}{4}\right) = -\frac{1}{4}\)，求 \(\cos 4x\) 的值.

\end{problem}

\begin{solution}
由积化和差公式，\(\frac{1}{2}\left[\sin(2x - \pi) + \sin\left(-\frac{\pi}{2}\right)\right] = -\frac{1}{4}\). 所以 \(\sin 2x = -\frac{1}{2}\). 从而 \(\cos 4x = 1 - 2\sin^2 2x = \frac{1}{2}\).
\end{solution}

\begin{problem}
设有一颗彗星，围绕地球沿一抛物线轨道运行，地球恰好位于这抛物线轨道的焦点处. 当此彗星离地球为 \(d\)（万公里）时，经过地球和彗星的直线与抛物线的轴的夹角为 \(30^{\circ}\). 求这彗星与地球的最短距离.

\end{problem}

\begin{solution}
【法一】 设极坐标系如图，使极点 \(O\) 位于抛物线的焦点处，\(Ox\) 过抛物线的轴. 则抛物线的极坐标方程为 \(\rho = \frac{p}{1 - \cos\theta}\). 由假设可得两种情形. 情形 1. 当 \(\theta = 30^{\circ}\) 时，\(\rho = d\)，故 \(d = \frac{p}{1 - \cos 30^{\circ}}\)，\(p = \left(1 - \frac{\sqrt{3}}{2}\right)d\). 当 \(\theta = \pi\) 时，\(\rho\) 有最小值 \(\rho_{\min} = \frac{p}{2} = \frac{2 - \sqrt{3}}{4}d\).

情形 2. 当 \(\theta = 150^{\circ}\) 时，\(\rho = d\)，故 \(d = \frac{p}{1 - \cos 150^{\circ}}\)，\(p = \left(1 + \frac{\sqrt{3}}{2}\right)d\). 当 \(\theta = \pi\) 时，\(\rho\) 有最小值 \(\rho_{\min} = \frac{p}{2} = \frac{2 + \sqrt{3}}{4}d\).

故这颗彗星与地球的最短距离为 \(\frac{2 - \sqrt{3}}{4}d\) 或 \(\frac{2 + \sqrt{3}}{4}d\)（万公里）.

【法二】 设平面直角坐标系 \(xOy\)，使 \(x\) 轴过抛物线的轴，抛物线方程为 \(y^2 = 2px\)（\(p > 0\)）. 按假设，地球位于 \(\left(\frac{p}{2},0\right)\) 处. 当彗星离地球为 \(d\) 时，彗星的位置为 \(M\left(\frac{p}{2} + d\cos 30^{\circ}, d\sin 30^{\circ}\right)\) 或 \(M'\left(\frac{p}{2} - d\cos 30^{\circ}, d\sin 30^{\circ}\right)\). 把 \(M\) 与 \(M'\) 的坐标分别代入 \(y^2 = 2px\)，得 \(p^2 \pm \sqrt{3}dp - \frac{d^2}{4} = 0\)，解得 \(p = \left(1 \mp \frac{\sqrt{3}}{2}\right)d\).

当彗星位于 \(P(x,y)\) 时，它与地球的距离为 \(\sqrt{\left(x - \frac{p}{2}\right)^2 + y^2} = x + \frac{p}{2} \geqslant \frac{p}{2}\)（\(x \geqslant 0\)）. 故当 \(x = 0\) 时取得最短距离，最短距离为 \(\frac{p}{2} = \frac{2 - \sqrt{3}}{4}d\) 或 \(\frac{p}{2} = \frac{2 + \sqrt{3}}{4}d\). 因此，这颗彗星与地球的最短距离为 \(\frac{2 - \sqrt{3}}{4}d\) 或 \(\frac{2 + \sqrt{3}}{4}d\)（万公里）.
\end{solution}

\begin{problem}
解关于实数 \(x\) 的不等式 \(\left|(\log_a x)^2 - 1\right| > 2a - 1\)，其中 \(a > 0\)，\(a \neq 1\).

\end{problem}

\begin{solution}
当 \(0 < a < \frac{1}{2}\) 时，原不等式的解为 \(x > 0\).

当 \(a = \frac{1}{2}\) 时，原不等式的解为 \(0 < x < \frac{1}{2}\)，\(\frac{1}{2} < x < 2\)，或 \(x > 2\).

当 \(a > \frac{1}{2}\) 时，分两种情形讨论.

\(1^{\odot}\) 若 \(\left(\log_a x\right)^2 - 1 > 2a - 1\)，则 \(\left(\log_a x\right)^2 > 2a\). 于是 \(\log_a x < -\sqrt{2a}\) 或 \(\log_a x > \sqrt{2a}\). 因此，当 \(\frac{1}{2} < a < 1\) 时，\(x > a^{-\sqrt{2a}}\) 或 \(0 < x < a^{\sqrt{2a}}\)；当 \(a > 1\) 时，\(0 < x < a^{-\sqrt{2a}}\) 或 \(x > a^{\sqrt{2a}}\).

\(2^{\odot}\) 若 \(-\left[\left(\log_a x\right)^2 - 1\right] > 2a - 1\)，则 \(\left(\log_a x\right)^2 < 2(1 - a)\). 这时必有 \(\frac{1}{2} < a < 1\)，且 \(-\sqrt{2(1 - a)} < \log_a x < \sqrt{2(1 - a)}\). 故 \(a^{\sqrt{2(1 - a)}} < x < a^{-\sqrt{2(1 - a)}}\).

综上，原不等式的解为
\[
\begin{cases}
x > 0, & 0 < a < \frac{1}{2}, \\[4pt]
0 < x < \frac{1}{2}\ \text{或}\ \frac{1}{2} < x < 2\ \text{或}\ x > 2, & a = \frac{1}{2}, \\[4pt]
0 < x < a^{\sqrt{2a}}\ \text{或}\ a^{\sqrt{2(1 - a)}} < x < a^{-\sqrt{2(1 - a)}}\ \text{或}\ x > a^{-\sqrt{2a}}, & \frac{1}{2} < a < 1, \\[4pt]
0 < x < a^{-\sqrt{2a}}\ \text{或}\ x > a^{\sqrt{2a}}, & a > 1.
\end{cases}
\]
\end{solution}

\begin{problem}
如图，设 \(\triangle ABC\) 和 \(\triangle DBC\) 所在的两平面互相垂直，且 \(AB = BC = BD\)，\(\angle CBA = \angle DBC = 120^{\circ}\). 求

\begin{enumerate}
\item \(A\)，\(D\) 的连线和平面 \(BCD\) 所成的角；
\item \(A\)，\(D\) 的连线和直线 \(BC\) 所成的角；
\item 二面角 \(A-BD-C\) 的大小（用反三角函数表示）.
\end{enumerate}

\bitmapfigure[max width=0.42\linewidth]{img/1991/shanghai/q24_fig1.png}

\end{problem}

\begin{solution}
连 \(A\)，\(D\). 过 \(A\) 点在平面 \(ABC\) 内作 \(BC\) 的垂线，交 \(CB\) 的延长线于 \(H\)，连 \(H\)，\(D\). 由面面垂直的性质，\(AH \perp\) 平面 \(BCD\)，而 \(HD\) 为 \(AD\) 在平面 \(BCD\) 内的射影，所以 \(\angle ADH\) 即为 \(AD\) 与平面 \(BCD\) 所成的角.
由题设可知 \(\triangle AHB \cong \triangle DHB\)，故 \(DH \perp HB\)，\(AH = DH\). 因而 \(\angle ADH = 45^{\circ}\). 所以 \(AD\) 与平面 \(BCD\) 所成的角为 \(45^{\circ}\).

又因为 \(HD\) 为 \(AD\) 在平面 \(BCD\) 内的射影，且 \(BC \perp HD\)，故 \(AD \perp BC\). 所以 \(A\)，\(D\) 的连线和 \(BC\) 所成的角为 \(90^{\circ}\).

在平面 \(BCD\) 内，过 \(H\) 作 \(HR \perp BD\)，垂足为 \(R\)，连 \(A\)，\(R\). 由三垂线定理，得 \(AR \perp BD\). 因而 \(\angle ARH\) 为二面角 \(A-BD-C\) 的平面角的补角.
设 \(BC = a\). 由题设得 \(AH = DH = \frac{\sqrt{3}}{2}a\)，\(HB = \frac{1}{2}a\). 在 \(\triangle HDB\) 内，可求得 \(HR = \frac{\sqrt{3}}{4}a\). 所以 \(\tan\angle ARH = \frac{AH}{HR} = 2\)，即 \(\angle ARH = \arctan 2\). 故二面角 \(A-BD-C\) 的大小为 \(\pi - \arctan 2\).
\end{solution}

\begin{problem}
\begin{enumerate}
\item 设复数 \(z\) 的辐角为 \(\theta\)（\(0 \leqslant \theta < \pi\)），且满足等式 \(|z - \i| = 1\)，求复数 \(z^2 - z\i\) 的辐角（用含 \(\theta\) 的式子表示），其中 \(\i\) 为虚数单位；
\item 设复数 \(z\) 满足等式 \(|z - \i| = 1\)，且 \(z \neq 0\)，\(z \neq 2\i\). 又复数 \(\omega\) 使得 \(\frac{\omega}{\omega - 2\i} \cdot \frac{z - 2\i}{z}\) 为实数. 问复数 \(\omega\) 在复平面上所对应的点 \(Z\) 的集合是什么图形，并说明理由.
\end{enumerate}

\end{problem}

\begin{solution}
【法一】 （1）设 \(z = r(\cos \theta + \i \sin \theta)\)（\(0 \leqslant \theta < \pi\)）. 由 \(|z - \i| = 1\) 得 \((r\cos \theta)^2 + (r\sin \theta - 1)^2 = 1\)，即 \(r^2 - 2r\sin \theta = 0\). 所以 \(r = 2\sin \theta\)（\(0 < \theta < \pi\)）. 于是 \(z^2 - z\i = z(z - \i) = 2\sin \theta \left[\cos\left(3\theta - \frac{\pi}{2}\right) + \i \sin\left(3\theta - \frac{\pi}{2}\right)\right]\). 故 \(z^2 - z\i\) 的辐角为 \(\left(3\theta - \frac{\pi}{2}\right) + 2k\pi\)（\(k \in \Z\)）.

（2）设 \(z = a + b\i\)，\(\omega = x + y\i\). 记 \(u = \frac{\omega}{\omega - 2\i} \cdot \frac{z - 2\i}{z}\). 由 \(|z - \i| = 1\) 可得 \(a^2 + b^2 - 2b = 0\). 又因为 \(u\) 为实数，所以 \(\operatorname{Im}u = 0\). 计算得 \(a(x^2 + y^2 - 2y) = 0\). 由于 \(z \neq 0\) 且 \(z \neq 2\i\)，故 \(a \neq 0\). 因而 \(x^2 + y^2 - 2y = 0\)，即 \(x^2 + (y - 1)^2 = 1\). 所以，复数 \(\omega\) 在复平面上所对应的点 \(Z\) 的集合是以 \((0,1)\) 为圆心，\(1\) 为半径的圆，除去点 \((0,2)\).

【法二】 设 \(z = 2\sin\theta(\cos\theta + \i\sin\theta)\)（\(0 < \theta < \frac{\pi}{2}\)，或 \(\frac{\pi}{2} < \theta < \pi\)），则 \(z - 2\i = 2\cos\theta\left[\cos\left(\theta - \frac{\pi}{2}\right) + \i\sin\left(\theta - \frac{\pi}{2}\right)\right]\)，
\(\frac{z - 2\i}{z} = \cot\theta\left[\cos\left(-\frac{\pi}{2}\right) + \i\sin\left(-\frac{\pi}{2}\right)\right] = -\i\cot\theta\).
\bitmapfigure[max width=0.52\linewidth]{img/1991/shanghai/q25_fig1.png}

又如图，\(\alpha = \frac{\pi}{2}\)，故 \(\arg\left(\frac{z - 2\i}{z}\right) = \frac{\pi}{2}\) 或 \(\frac{3\pi}{2}\). 按照复数的乘法，除法的几何意义与假设条件 \(\operatorname{Im}u = 0\)，又可得 \(\arg\left(\frac{\omega - 0}{\omega - 2\i}\right) + \arg\left(\frac{z - 2\i}{z}\right) = k\pi\)（\(k = 1\)，\(2\)，\(3\)）.
故 \(\arg\left(\frac{\omega - 0}{\omega - 2\i}\right) = \frac{\pi}{2}\) 或 \(\frac{3\pi}{2}\). 从而 \(\beta = \frac{\pi}{2}\). 这说明点 \(Z\) 必在圆 \(|\omega - \i| = 1\) 上. 因此 \(\omega\) 在复平面上所对应的点 \(Z\) 的集合是以 \((0,1)\) 为圆心，\(1\) 为半径的圆，除去点 \((0,2)\).
\end{solution}

\begin{problem}
如图，设有一动直线 \(L\)，过定点 \(A(2,0)\) 且与抛物线 \(y = x^2 + 2\) 相交于不同的两点 \(B\) 和 \(C\). 点 \(B\)，\(C\) 在 \(x\) 轴上的射影分别是 \(B'\)，\(C'\). \(P\) 是线段 \(BC\) 上的点，适合关系式 \(\frac{BP}{PC} = \frac{|BB'|}{|CC'|}\). 求 \(\triangle POA\) 的重心 \(Q\) 的轨迹方程，并说明轨迹是什么图形.

\bitmapfigure[max width=0.42\linewidth]{img/1991/shanghai/q26_fig1.png}

\end{problem}

\begin{solution}
设点 \(P(x_0,y_0)\)，\(B(x_1,y_1)\)，\(C(x_2,y_2)\)，\(\triangle POA\) 的重心为 \(Q(X,Y)\). 又设 \(\frac{y_1}{y_2} = \lambda\). 由题意，\(y_0 = \frac{y_1 + \lambda y_2}{1 + \lambda} = \frac{2y_1y_2}{y_1 + y_2}\). 可设直线 \(L\) 的方程为 \(y = k(x - 2)\)（\(k \neq 0\)）. 由此得 \(x = \frac{y + 2k}{k}\). 代入抛物线 \(y = x^2 + 2\)，得 \(y^2 + (4k - k^2)y + 6k^2 = 0\). 因此 \(y_1 + y_2 = k^2 - 4k\)，\(y_1y_2 = 6k^2\). 代入前式，得 \(\frac{2}{y_0} = \frac{y_1 + y_2}{y_1y_2} = \frac{k - 4}{6k}\). 又由 \(k = \frac{y_0}{x_0 - 2}\)，得 \(4x_0 - y_0 + 4 = 0\). 因为 \(Q\) 是 \(\triangle POA\) 的重心，所以 \(X = \frac{x_0 + 2}{3}\)，\(Y = \frac{y_0}{3}\). 由此得 \(Q\) 满足 \(12X - 3Y - 4 = 0\).

又因直线 \(L\) 与抛物线有两个相异实根，故判别式满足 \((4k - k^2)^2 - 24k^2 > 0\)，即 \(k < 4 - 2\sqrt{6}\) 或 \(k > 4 + 2\sqrt{6}\). 化到 \(Y\) 上，可得 \(4 - \frac{4}{3}\sqrt{6} < Y < 4 + \frac{4}{3}\sqrt{6}\)，且点 \(\left(\frac{4}{3},4\right)\) 不在轨迹上.

因此，\(\triangle POA\) 的重心 \(Q\) 的轨迹是直线 \(12X - 3Y - 4 = 0\) 上满足 \(4 - \frac{4}{3}\sqrt{6} < Y < 4 + \frac{4}{3}\sqrt{6}\) 的一段，并除去点 \(\left(\frac{4}{3},4\right)\).
\end{solution}

\begin{problem}
如图，在 \(\triangle ABC\) 中，\(BC = a\)，\(AC = b\)，\(AB = c\)，\(\angle ACB = \theta\). 现将 \(\triangle ABC\) 分别以 \(BC\)，\(AC\)，\(AB\) 所在的直线为轴旋转一周，设所得三个旋转体的体积依次为 \(V_1\)，\(V_2\)，\(V_3\). 求

\begin{enumerate}
\item \(T = \frac{V_3}{V_1 + V_2}\)（用 \(a\)，\(b\)，\(c\)，\(\theta\) 表示）；
\item 若 \(\theta\) 为定值，并令 \(\frac{a + b}{c} = x\)，将 \(T\) 表为 \(x\) 的函数，写出这函数的定义域，并求这函数的最大值 \(u\)；
\item 当 \(\theta\) 在 \(\left[\frac{\pi}{3},\pi\right)\) 内变化时，求 \(u\) 的最大值.
\end{enumerate}

\bitmapfigure[max width=0.38\linewidth]{img/1991/shanghai/q27_fig1.png}

\end{problem}

\begin{solution}
在 \(\triangle ABC\) 中，设分别以 \(BC\)，\(AC\)，\(AB\) 为底边的高依次为 \(h_1\)，\(h_2\)，\(h_3\). 则 \(h_1 = b\sin \theta\)，\(h_2 = a\sin \theta\)，\(h_3 = \frac{ab\sin \theta}{c}\). 于是 \(V_1 = \frac{1}{3}\pi h_1^2 a = \frac{\pi}{3}ab^2\sin^2 \theta\)，\(V_2 = \frac{1}{3}\pi h_2^2 b = \frac{\pi}{3}a^2b\sin^2 \theta\)，\(V_3 = \frac{1}{3}\pi h_3^2 c = \frac{\pi}{3}\frac{a^2b^2}{c}\sin^2 \theta\). 因此 \(T = \frac{V_3}{V_1 + V_2} = \frac{ab}{(a + b)c}\).

由余弦定理，\(c^2 = a^2 + b^2 - 2ab\cos \theta = (a + b)^2 - 2ab(1 + \cos \theta)\). 令 \(x = \frac{a + b}{c}\)，则 \(ab = \frac{(x^2 - 1)c^2}{2(1 + \cos \theta)}\). 代入上式，得 \(T = \frac{1}{2(1 + \cos \theta)}\left(x - \frac{1}{x}\right)\).
又因为 \(x = \frac{a + b}{c} > 1\)，且 \(c^2 = (a + b)^2 - 2ab(1 + \cos \theta) \geqslant \frac{1 - \cos \theta}{2}(a + b)^2\)，所以 \(x \leqslant \sqrt{\frac{2}{1 - \cos \theta}}\)，当且仅当 \(a = b\) 时取等号. 因此函数的定义域为 \(\left(1,\sqrt{\frac{2}{1 - \cos \theta}}\right]\).
因为 \(T\) 在该区间内单调递增，所以当 \(x = \sqrt{\frac{2}{1 - \cos \theta}}\)，即 \(a = b\) 时，\(T\) 取得最大值 \(u = \frac{\sqrt{2}}{4\sqrt{1 - \cos \theta}}\).

当 \(\theta \in \left[\frac{\pi}{3},\pi\right)\) 时，\(\sqrt{1 - \cos \theta}\) 单调递增，从而 \(u\) 单调递减. 因此当 \(\theta = \frac{\pi}{3}\) 时，\(u\) 取得最大值 \(u_{\max} = \frac{1}{2}\).
\end{solution}
